\documentclass[12pt, a4paper]{article}

% ---- Standard Packages ----
\usepackage[a4paper, margin=1in]{geometry}
\usepackage[utf8]{inputenc}
\usepackage[T1]{fontenc}
\usepackage[bahasa]{babel}
\usepackage{amsmath, amssymb, amsthm, mathtools}
\usepackage{enumitem}
\usepackage{booktabs}
\usepackage{multicol}
\usepackage{hyperref}

% ---- Basic Metadata ----
\title{MODUL UTAMA PANGGAJARAN BASA SUNDA\\ KELAS 10 -- KOMPETENSI KEAHLIAN TKJ\\ \large Buku Panduan Belajar Mandiri Komprehensif \& Teoretis}
\author{Tim Pengajar Kompetensi Bahasa \& Sastra Sunda}
\date{}

\begin{document}

\maketitle

\begin{center}
    \textbf{Topik Utama Paparan Materi:} \\
    Komunikasi Wawancara $\cdot$ Historiografi Carita Babad $\cdot$ Silsilah Pancakaki $\cdot$ Apresiasi Sajak Modern
\end{center}

\tableofcontents
\newpage

% ==============================================================
\section{METODOLOGI JEUNG ETIKA WAWANCARA}
% ==============================================================

\subsection{Definisi jeung Esensi Komunikasi Wawancara}
Dina dunya profésional jeung komunikasi komunikasi massa, néangan informasi anu akurat merlukeun métode anu sistematis. Sacara harti dasar, \textbf{tanya jawab langsung pikeun nyajaring informasi ti narasumber disebut na wawancara}. Ieu kagiatan téh mangrupa prosés interaksi verbal anu boga tujuan pikeun ngumpulkeun data, fakta, atawa opini ngeunaan hiji pasualan.

Dina kagiatan wawancara, aya dua pihak anu silih boga lalakon:
\begin{enumerate}
    \item \textbf{Pewawancara (Nu Ngawawancara):} Jalma anu boga pancén pikeun ngadali prosés tanya jawab sarta ngajukeun patanyaan.
    \item \textbf{Narasumber (Informan):} Jalma anu dianggap miboga kaahlian, otoritas, atawa pangalaman langsung ngeunaan topik anu nuju ditalungtik.
\end{enumerate}

\subsection{Klasifikasi Wawancara Dumasar Perencanaan}
Dumasar kana cara nyiapkeun instruksina sarta kabébasan dina ngajukeun patanyaan, wawancara kabagi jadi sababaraha jinis:
\begin{itemize}
    \item \textbf{Wawancara Terstruktur:} Wawancara anu daptar pertanyaanana geus dirumuskeun sacara rapih jeung kaku sacara tinulis saméméh acara lumangsung.
    \item \textbf{Wawancara Bebas:} Nyaéta \textbf{kagiatan wawancara anu pananya teu dirumuskeun} sacara tinulis sateuacanna. Nu ngawawancara langsung mekarkeun patanyaan sacara spontan dumasar kana kaayaan jeung jawaban nu kaluar ti narasumber.
    \item \textbf{Wawancara Bebas Terpimpin:} Gabungan tina duanana, dimana nu ngawawancara mawa gurat badag (outline) topik tapi patanyaanana dimekarkeun sacara bébas di lapangan.
\end{itemize}

\subsection{Materi Tambahan: Rumus Kalimah Pananya (5W + 1H)}
Pikeun meunangkeun informasi anu jentre, saurang pewawancara kudu maham kana fungsi kecap pananya dina basa Sunda:
\begin{itemize}
    \item \textbf{Naon (What):} Pikeun nanyakeun kajadian, barang, atawa hal.
    \item \textbf{Saha (Who):} Pikeun nanyakeun tokoh, jalma, atawa palaku.
    \item \textbf{Mana / Di mana (Where):} Pikeun nanyakeun tempat atawa lokasi kajadian.
    \item \textbf{Iraha (When):} Pikeun nanyakeun waktu lumangsungna peristiwa.
    \item \textbf{Naha (Why):} Pikeun nanyakeun alesan atawa kasang tukang (latar belakang).
    \item \textbf{Kumaha / Sabaraha (How / How much):} Pikeun nanyakeun proses, kaayaan, atawa jumlah.
\end{itemize}

\subsection{Materi Tambahan: Tatakrama Basa dina Wawancara}
Saurang siswa TKJ anu rék ngalakukeun wawancara kudu niténan \textbf{Tatakrama Basa Sunda} (Ragaw Basa) salaku wujud sopan santun:
\begin{enumerate}
    \item **Basa Hormat (Lemes):** Wajib dipaké nalika ngawawancara jalma anu leuwih kolot, dipihormat, atawa nu anyar pinanggih. Contona ngagunakeun kecap \textit{gandeng} diganti ku \textit{bilih}, \textit{nanya} diganti ku \textit{naroskeun}.
    \item **Basa Loma (Akrab):** Ukur bisa dipaké upami narasumberna babaturan saumuran anu geus akrab pisan.
\end{enumerate}

\newpage

% ==============================================================
\section{HISTORIOGRAFI SASTRA: CARITA BABAD}
% ==============================================================

\subsection{Hakékat jeung Definisi Carita Babad}
Dina wangun prosa heubeul (buhun), tatar Sunda miboga karya sastra anu unik sabab nyampurkeun unsur fiksi jeung realitas sajarah lokal. \textbf{Carita atawa dongeng anu ngandung unsur sajarah lokal disebut babad}. Kalayan jentre, \textbf{carita babad teh kaasup kana wangun carita buhun anu eusina ngenaan kajadian atawa lalakon di jaman baheula}. 

\subsection{Karakteristik Mandiri jeung Ciri-Ciri Utama}
Carita babad miboga bédana anu kacida jentréna upami dibandingkeun jeung dongéng umum saperti sage atawa mitos. 
\begin{itemize}
    \item \textbf{Ciri Khas:} Nu jadi \textbf{ciri khas tina carita babad nyaeta miboga unsur sajarah}. Ieu hal bisa ditingali tina ayana rujukan ngaran tempat (toponimi), ngaran raja, sarta wilayah karajaan anu cocog jeung data sajarah sabenerna.
    \item \textbf{Unsur Pamohalan:} Sanajan aya unsur sajarahna, \textbf{dina carita babad, biasana aya unsur du dicampur antara sajarah jeung pamohalan}. Kecap *pamohalan* hartina hal-hal anu mustahil atawa teu asup akal (magis/fiktif). Contona: Tokoh utama miboga kasaktian anu bisa ngaleungit, merangan musuh ku kakuatan angin, atawa lahir tina awi.
\end{itemize}

Sacara sintésis, \textbf{ciri ciri carita babad (lalaokn baheula, pamohalan, miboga unsur sajarah)} mangrupa tilu tihang utama anu ngawangun ieu karya sastra buhun.

\subsection{Tujuan sarta Struktur Konstruksi (Galur)}
\begin{itemize}
    \item \textbf{Tujuan Penulisan:} Dina fungsi sosialna, \textbf{tujuan carita babad nyaeta nyalamat sajarah baheula} sangkan jati diri hiji komunitas, silsilah para raja, sarta kajayaan karajaan di jaman baheula henteu leungit sarta bisa dipikanyaho ku katurunan saterusna.
    \item \textbf{Struktur Narasi:} Runtuyan kajadian dina babad dikonstruksikeun ku plot anu rapih. \textbf{Runtuyan pirang pirang kajadian nu ngawangun dina carita babad disebut galur}. Galur anu pangseringna dipaké nyaéta galur maju (merélé dumasar urutan waktu).
\end{itemize}

\subsection{Materi Tambahan: Unsur Intrinsik Carita Babad}
Siga karya prosa séjénna, carita babad diwangun ku unsur intrinsik di antarana:
\begin{enumerate}
    \item \textbf{Téma:} Gagasan poko atawa ide pusat anu ngajiwaan sakabéh eusi carita.
    \item \textbf{Palaku (Tokoh):} Ngaran palaku anu aya dina carita, biasana dibagi jadi tokoh utama (raja/pahlawan) jeung tokoh pangrojong.
    \item \textbf{Latar (Setting):} Tempat (karaton, leuweung, nagara), waktu (jaman baheula), jeung suasana (perang, sedih, agung).
\end{enumerate}

\subsection{Karajaan Munggaran di Tatar Sunda jeung Contoh Babad}
Dina dimensi sajarah nyata, \textbf{karajaan munggaran di jawa barat nyaeta tarumanagara}. Karajaan Hindu aliran Waisnawa ieu ngadeg ti abad ka-4 Masehi sarta dipingpin ku raja kawentar nyaéta Purnawarman.

\vspace{0.2cm}
\noindent \textbf{Contoh-Contoh Carita Babad nu Kakoncara di Tatar Sunda:}
\begin{multicols}{2}
\begin{itemize}
    \item \textit{Babad Bogor:} Sajarah runtagna Karajaan Pajajaran.
    \item \textit{Babad Galuh:} Lalakon Karajaan Galuh di Ciamis.
    \item \textit{Babad Cirebon:} Mimiti ngadegna Kasultanan Islam Cirebon.
    \item \textit{Babad Banten:} Sajarah lokal Islamisasi di Banten.
    \item \textit{Babad Sumedang:} Lalakon sajarah Karajaan Sumedang Larang.
    \item \textit{Babad Limbangan:} Asal-usul sajarah daerah Garut.
\end{itemize}
\end{multicols}

\newpage

% ==============================================================
\section{SISTEM KEKERABATAN KULAWARGA: PANCAKAKI}
% ==============================================================

\subsection{Definisi jeung Fungsi Sosial Pancakaki}
Dina budaya Sunda, papatokan pikeun nataan hubungan kekeluargaan kacida dijagana. \textbf{Naon anu dimaksud pancakaki teh (silsilah kaluwarga)} atawa tatanan anu nuduhkeun perenahna duduluran atawa kekerabatan ti hiji jalma ka jalma lianna. Fungsi utamana nyaéta pikeun ngajaga tali silaturahmi sarta nerapkeun tatakrama basa anu merenah dumasar tingkatan umur dunya kulawarga.

\subsection{Rundayan Pancakaki ka Luhur (Ascending)}
Nataan silsilah ka luhur hartina nataan karuhun urang (garis katurunan saluhureun kolot urang):

\begin{center}
    \small
    \begin{tabular}{clp{8.5cm}}
        \toprule
        \textbf{No.} & \textbf{Istilah Katurunan} & \textbf{Harti Perenahna Hubungan Kulawarga} \\
        \midrule
        1 & Kolot & Bapa atawa indung kandung urang. \\
        2 & Aki / Nini & Kolotna ti bapa atawa indung urang. \\
        3 & Buyut & Kolotna ti aki atawa nini urang. \\
        4 & \textbf{Bao} & \textbf{Kolot na buyut sok disebut bao}. \\
        5 & Janggawareng & Kolotna ti bao urang. \\
        6 & Udeg-udeg & Kolotna ti janggawareng urang. \\
        7 & \textbf{Gantung Siwur} & \textbf{Gantung siwur dina pancakaki nyaeta bapana udeg udeg}. \\
        \bottomrule
    \end{tabular}
\end{center}

\vspace{0.3cm}
\subsection{Materi Tambahan: Rundayan Pancakaki ka Handap (Descending)}
Salian ti nataan ka luhur, urang ogé kudu nyaho nataan rundayan katurunan ka handap (anak-incu urang ka hareupna):

\begin{center}
    \small
    \begin{tabular}{clp{8.5cm}}
        \toprule
        \textbf{No.} & \textbf{Istilah Katurunan} & \textbf{Harti Perenahna Hubungan Kulawarga} \\
        \midrule
        1 & Anak & Katurunan langsung kahiji ti urang sarta salaki/pamajikan. \\
        2 & Incu & Anakna ti anak urang. \\
        3 & Umpi & Anakna ti incu urang. \\
        4 & Cicip & Anakna ti umpi urang. \\
        5 & Munyang & Anakna ti cicip urang. \\
        6 & Gewari & Anakna ti munyang urang. \\
        7 & Saompi & Anakna ti gewari urang. \\
        \bottomrule
    \end{tabular}
\end{center}

\vspace{0.3cm}
\subsection{Istilah Hubungan Horizontal jeung Diagonal (Duduluran)}
Dina tatanan gigir (samping), aya istilah husus anu nuduhkeun darajat kadekatan getih:
\begin{itemize}
    \item \textbf{Dulur Pet Ku Hinis:} \textbf{Istilah dulur pet kuhinis dina pancakaki hartina dulur kandung}, nyaéta dulur anu kaluar tina rahim indung sarta bapa anu sarua.
    \item \textbf{Bau-Bau Sinduk:} \textbf{Istilah bau bau sinduk dina pancakaki hartina baraya keneh sanajan tos laen}. Ieu nuduhkeun yén hubungan duduluranana geus jauh pisan (contona dulur misan tingkat kaopat atawa kalima), tapi sacara adat masih dianggap sakulawarga.
    \item \textbf{Alo jeung Suwar:} Hubungan jeung anak ti dulur urang miboga patokan: lamun urang \textbf{kaanak lanceuk disebut alo, ari ka anak adi disebut suwar}.
\end{itemize}

\newpage

% ==============================================================
\section{ANALISIS PUISI MODERN: SAJAK}
% ==============================================================

\subsection{Definisi jeung Sejarah Kamekaran Sajak}
Sajak mangrupa salah sahiji wangun puisi modern anu aya dina sastra Sunda. \textbf{Naon anu dimaksud sajak teh (puisi)} wangun bébas anu henteu kabeungkeut ku aturan-aturan ketat saperti pupuh (anu boga guru lagu jeung guru wilangan). Sajak asup kana sastra Sunda dina jaman sabada merdeka (sabudeureun taun 1946), dilantarankeun ku pangaruh panyajak Indonesia.

\subsection{Struktur Lahir (Unsur Intrinsik) dina Sajak}
Pikeun neuleuman hiji sajak, urang kudu maham kana elemen panyusun strukturna:
\begin{itemize}
    \item \textbf{Padalisan:} Dina wangunan fisik atanapi bait, \textbf{naon ari anu disebut padalisan teh (bait sajak)} atawa jajaran kalimat anu ngawangun hiji harti dina sajak.
    \item \textbf{Amanat:} \textbf{Nu disebut amanan dina sajak teh hartina pesan moral} atawa talatah, harepan, jeung nasehat anu hayang ditepikeun ku panyajak ka nu maca liwat karyana.
    \item \textbf{Imaji:} \textbf{Naon ari anu disebut imaji teh} nyaéta \textbf{citraan atawa gambaran angen-angen/pikiran} anu bisa ngahudangkeun panca indra (visual, auditif, taktil) sangkan nu maca saolah-olah ngadéngé, nempo, atawa ngarasa sorangan naon anu dicaritakeun dina sajak.
\end{itemize}

\subsection{Gaya Basa jeung Tujuan Estetis}
\begin{itemize}
    \item \textbf{Pilihan Kecap (Diksi):} \textbf{Gaya bahasa anu sok dipake dina sajak biasana bahasa endah}. Basa anu dipaké nyaéta basa anu pinuh ku simbol, kiasan (metafora), sarta miboga harti konotatif (lain harti sabenerna) sangkan éndah dirasakeunana.
    \item \textbf{Tujuan Penulisan:} Unggal seniman boga udagan séwang-séwangan, tapi sacara umum \textbf{sebutkeun salasahiji tujuan nyieun sajak jang nyatakeun rasa, pikiran jeung pangalaman} batin, éksprési jiwa, sarta salaku média kritik sosial kaayaan sabudeureun hirupna.
\end{itemize}

\subsection{Materi Tambahan: Struktur Batin Sajak}
Salian ti struktur lahir di luhur, sajak ogé diwangun ku \textbf{Struktur Batin} di antarana:
\begin{enumerate}
    \item \textbf{Rasa (Feeling):} Sikep panyajak kana masalah utama anu aya dina sajakna (contona: kaayaan sedih, ambek, keueung, atawa gumbira).
    \item \textbf{Nada (Tone):} Sikep panyajak ka nu maca (contona: nada mapatahan, nyindir, atawa ngajak babarengan).
\end{enumerate}

\subsection{Langkah Strategis Nulis jeung Tokoh Sajak Sunda}
\begin{itemize}
    \item \textbf{Tokoh Panyajak Kasohor:} \textbf{Saha ngaran panyajak anu kasohor, jeung dianggap salasahiji seniman sajak sastra sunda nu nyebarkeun sajak (kis ws)}. Ngaran lengkepna nyaéta Kiswanda Sastradinata. Anjeunna mangrupa pelopor utama anu ngajadikeun sajak ditarima ku masyarakat Sunda dina mangsa awal kamekaranana.
    \item \textbf{Langkah Nulis Sajak:} \textbf{Samemeh rek nulis sajak aya sabaraha hal anu kudu di cangkumm (6)} perkara utama anu jadi sarat sangkan sajakna alus, nyaéta:
    \begin{enumerate}
        \item Nangtukeun téma (ide dasar sajak).
        \item Nangtukeun suasana atawa milarian inspirasi.
        \item Milih diksi (pamilihan kecap anu éndah sarta merenah).
        \item Ngolah imaji (nyiptakeun citraan panca indra).
        \item Ngatur wirahma (ritme jeung musikalitas sora tungtung).
        \item Nangtukeun amanat (pesen moral anu hayang ditepikeun).
    \end{enumerate}
\end{itemize}

\end{document}
